\diarytitle{0078}{行列方程式 AX=B の解（逆行列を用いた解法）}

$A$ が正則ならば、$AX = B$ の両辺に左から $A^{-1}$ を掛けて $X = A^{-1}B$ と一意に解ける。まず $A$ が正則であることを確かめ、$A^{-1}$ を求める。

第 1 行に関して余因子展開すると
\[
\det A = 1\cdot\begin{vmatrix}1&1\\0&1\end{vmatrix} - 1\cdot\begin{vmatrix}0&1\\1&1\end{vmatrix} + 0
= 1\cdot 1 - 1\cdot(-1) = 2 \neq 0
\]
であるから $A$ は正則である。余因子を計算すると
\[
A_{11}=1,\ A_{12}=1,\ A_{13}=-1,\quad
A_{21}=-1,\ A_{22}=1,\ A_{23}=1,\quad
A_{31}=1,\ A_{32}=-1,\ A_{33}=1
\]
より随伴行列は
\[
\operatorname{adj}A =
\begin{pmatrix}
1 & -1 & 1 \\
1 & 1 & -1 \\
-1 & 1 & 1
\end{pmatrix},
\qquad
A^{-1} = \frac{1}{2}
\begin{pmatrix}
1 & -1 & 1 \\
1 & 1 & -1 \\
-1 & 1 & 1
\end{pmatrix}
\]
$A$ は正則だから $X = A^{-1}B$ と一意に定まる。$A^{-1}$ の各行と $B$ の各列の内積を計算すると
\[
X = \frac{1}{2}
\begin{pmatrix}
1 & -1 & 1 \\
1 & 1 & -1 \\
-1 & 1 & 1
\end{pmatrix}
\begin{pmatrix}
4 & 3 \\
3 & 5 \\
1 & 6
\end{pmatrix}
= \frac{1}{2}
\begin{pmatrix}
4-3+1 & 3-5+6 \\
4+3-1 & 3+5-6 \\
-4+3+1 & -3+5+6
\end{pmatrix}
= \frac{1}{2}
\begin{pmatrix}
2 & 4 \\
6 & 2 \\
0 & 8
\end{pmatrix}
\]
したがって
\[
\boxed{\; X =
\begin{pmatrix}
1 & 2 \\
3 & 1 \\
0 & 4
\end{pmatrix} \;}
\]
検算として $AX$ を直接計算する。第 1 行 $(1,1,0)$ と $X$ の各列の内積は $(1+3,\,2+1)=(4,3)$、第 2 行 $(0,1,1)$ とは $(3+0,\,1+4)=(3,5)$、第 3 行 $(1,0,1)$ とは $(1+0,\,2+4)=(1,6)$ となる。すなわち
\[
AX =
\begin{pmatrix}
4 & 3 \\
3 & 5 \\
1 & 6
\end{pmatrix} = B
\]
であり、$B$ に一致することが確かめられる。
